% Ad Familiares 7.17 (ad-familiares-07-17)
% Source: https://raw.githubusercontent.com/PerseusDL/canonical-latinLit/master/data/phi0474/phi056/phi0474.phi056.perseus-lat1.xml
% Fetched by scripts/fetch_latin.py

% Dateline (Perseus): Scr. Romae ex. m. Od. a. 700 (54) .

\ciceroLetterOpener{CICERO TREBATIO S.}

\ciceroSection{1} ex tuis litteris et Quinto fratri gratias egi et te aliquando conlaudare possum, quod iam videris certa aliqua in sententia constitisse. nam primorum mensum litteris tuis vehementer commovebar, quod mihi interdum (pace tua dixerim) levis in urbis urbanitatisque desiderio, interdum piger, interdum timidus in labore militari, saepe autem etiam, quod a te alienissimum est, subimpudens videbare. tamquam enim syngrapham ad imperatorem non epistulam attulisses, sic pecunia ablata domum redire properabas, nec tibi in mentem veniebat eos ipsos qui cum syngraphis venissent Alexandream nummum adhuc nullum auferre potuisse.

\ciceroSection{2} ego si mei commodi rationem ducerem, te mecum esse maxime vellem; non enim mediocri adficiebar vel voluptate ex consuetudine nostra vel utilitate ex consilio atque opera tua; sed cum te ex adulescentia tua in amicitiam et fidem meam contulisses; semper te non modo tuendum mihi sed etiam augendum atque ornandum putavi. itaque quoad opinatus sum me in provinciam exiturum, quae ad te ultro detulerim meminisse te credo. postea quam ea mutata ratio est, cum viderem me a Caesare honorificentissime tractari et unice diligi hominisque liberalitatem incredibilem et singularem fidem nossem, sic ei te commendavi et tradidi, ut gravissime diligentissimeque potui. quod ille ita et accepit et mihi saepe litteris significavit et tibi et verbis et re ostendit mea commendatione sese valde esse commotum. hunc tu virum nactus, si me aut sapere aliquid aut velle tua causa putas, ne dimiseris et, si quae te forte res aliquando offenderit, cum ille aut occupatione aut difficultate tardior tibi erit visus, perferto et ultima exspectato; quae ego tibi iucunda et honesta praestabo.

\ciceroSection{3} pluribus te hortari non debeo; tantum moneo, neque amicitiae confirmandae clarissimi ac liberalissimi viri neque uberioris provinciae neque aetatis magis idoneum tempus, si hoc amiseris, te esse ullum umquam reperturum. ' hoc ,' quem ad modum vos scribere soletis in vestris libris, 'idem Q. Cornelio videbatur.' in Britanniam te profectum non esse gaudeo, quod et labore caruisti et ego te de rebus illis non audiam. Ubi sis hibernaturus et qua spe aut condicione perscribas ad me velim.
